\addtocontents{toc}{\protect\Needspace{3\baselineskip}}
\originalchapter{简并星、吸积与恒星形成}\label{ch:compact}
\originalsection{电子简并与多方球}
白矮星不依靠高温维持压力. Pauli 原理使电子必须占据不同动量态. 零温时填满 $p\le p_F$, 两自旋态给 $n_e=p_F^3/(3\pi^2\hbar^3)$. 各向同性动量流的压强为
\begin{equation}
P=\frac1{3\pi^2\hbar^3}\int_0^{p_F}p^3v(p)\dd p.
\end{equation}
非相对论时 $v=p/m_e$, 因而 $P=K_{\rm NR}\rho^{5/3}$; 极相对论时 $v=c$, 因而 $P=K_{\rm R}\rho^{4/3}$. 其中 $\rho=\mu_em_pn_e$, $\mu_e$ 为每电子对应的质子质量倍数, 并非恒星气体平均分子量 $\mu$.

令 $P=K\rho^{1+1/n}$、$\rho=\rho_c\theta^n$, 取
$r=\alpha\xi$、$\alpha^2=(n+1)K\rho_c^{1/n-1}/(4\pi G)$.
把静水平衡式乘 $r^2/\rho$ 再求导, 使用质量守恒消去 $m_r$, 得
\begin{equation}\label{eq:lane}
\frac1{\xi^2}\frac{\dd}{\dd\xi}\left(\xi^2\theta'\right)=-\theta^n,
\qquad\theta(0)=1,\quad\theta'(0)=0.
\end{equation}
中心看似奇异, 正则展开却给 $\theta=1-\xi^2/6+O(\xi^4)$, 可由小正 $\xi$ 启动数值积分. 第一零点 $\xi_1$ 定义表面, $R=\alpha\xi_1$、$M=4\pi\alpha^3\rho_c[-\xi_1^2\theta'(\xi_1)]$. $n=3/2$ 时 $R\propto M^{-1/3}$; $n=3$ 时质量与 $\rho_c$ 无关, 给出 Chandrasekhar 质量尺度. 真实白矮星还受有限温度、库仑修正和逆 $\beta$ 过程影响.

\originalsection{吸积盘}
对质量 $M$ 的中心天体, 半径 $r$ 的 Kepler 速度为 $v_K=(GM/r)^{1/2}$. 圆轨道单位质量总能为 $-GM/(2r)$, 进入更小半径必须失去能量和角动量. 因此吸积不仅是物质向内掉落, 还需要应力把角动量向外输运. 薄盘中垂直重力近似 $-\Omega_K^2z$, 静水平衡给尺度高 $H\sim c_s/\Omega_K$; $H/r\ll1$ 要求 $c_s\ll v_K$.

对稳态、Kepler 薄盘, 以 $\dot M>0$ 为向内质量流率, 角动量守恒使应力矩满足 $\mathcal G=\dot M(\sqrt{GMr}-\sqrt{GMr_{\rm in}})$, 此处使用内缘零应力条件. 每环机械功率来自相邻环的相对角速度, 两面各自辐射通量为
\begin{equation}
F(r)=\frac{3GM\dot M}{8\pi r^3}\left(1-\sqrt{\frac{r_{\rm in}}r}\right).
\end{equation}
积分两面得到 $L=GM\dot M/(2r_{\rm in})$. 这个结果不适用于厚盘、强风、内缘有应力或强相对论区域. 辐射作用于自由电子, 通过电荷耦合带动离子; 比较 $\sigma_TL/(4\pi r^2c)$ 与质子引力 $GMm_p/r^2$ 得 $L_{\rm Edd}=4\pi GMm_pc/\sigma_T$. 这是球对称散射模型的力平衡尺度, 不表示所有盘必须严格低于它.

\originalsection{云的引力不稳定}
均匀气体中取小扰动 $\delta\rho,\boldsymbol v,\delta\Phi\propto\ee^{\ii(\boldsymbol k\cdot\boldsymbol x-\omega t)}$. 连续方程给 $-\omega\delta\rho+\rho_0\boldsymbol k\cdot\boldsymbol v=0$, Euler 方程给 $\omega\boldsymbol v=\boldsymbol k(c_s^2\delta\rho/\rho_0+\delta\Phi)$, Poisson 方程给 $-k^2\delta\Phi=4\pi G\delta\rho$. 消去两未知量得 $\omega^2=c_s^2k^2-4\pi G\rho_0$. 当 $k<k_J=\sqrt{4\pi G\rho_0}/c_s$ 时扰动增长. 本卷定义 $\lambda_J=2\pi/k_J$, Jeans 半径为 $R_J=\lambda_J/2$, 故 $M_J=\pi^{5/2}c_s^3/(6G^{3/2}\rho_0^{1/2})$. 不同文献的 Jeans 质量常数可能不同, 必须先声明约定.

这项无限均匀介质线性化以减去背景引力为理想化, 不是有限球的精确边界值解. 坍缩后若近似等温, $P$ 项相对于重力会越来越弱; 角动量和磁场却可阻止简单球对称坍缩, 所以形成盘、碎裂和角动量输运是恒星形成不可忽略的后续步骤.

\originalsection{Lane--Emden 方程的检验例}
从静水平衡可写 $\frac1{r^2}\frac{\dd}{\dd r}(\frac{r^2}{\rho}\frac{\dd P}{\dd r})=-4\pi G\rho$. 代 $P=K\rho_c^{1+1/n}\theta^{n+1}$、$\rho=\rho_c\theta^n$, 内部导数的 $\theta^n$ 与 $\rho$ 消去, 尺度 $\alpha$ 则消去剩余常数, 得式\eqref{eq:lane}. 此推导要求 $n>0$ 与正密度; $n=0$ 的不可压均匀密度球应直接取相应极限方程.

对 $n=0$, 直接积分 $\xi^2\theta'=-\xi^3/3$, 得 $\theta=1-\xi^2/6$, 第一零点 $\sqrt6$. 对 $n=1$, 令 $u=\xi\theta$, 方程化为 $u''+u=0$, 正则初值 $u(0)=0,u'(0)=1$, 故 $\theta=\sin\xi/\xi$, 第一零点 $\pi$, 质量参数 $-\pi^2\theta'(\pi)=\pi$. 两个解析例能检验中心处理与数值程序的归一化, 不能因它们可积就假定 $n=3/2,3$ 也有相同闭式.

无量纲曲线的作用是把中心密度与结构形状分开. 对固定 $n,K$, 改变 $\rho_c$ 只重缩放 $\alpha$ 和总质量, 无量纲 $\xi_1,w_n$ 不变. 因此质量半径尺度律是整族模型的结论, 不是某一次中心密度计算的巧合. 接近相对论简并极限时仍要注意, 真正的状态方程从 $5/3$ 连续过渡到 $4/3$, 并不是突然切换一个全星固定多方指数.
